\documentclass[11pt]{article}
\pagestyle{plain}

% \pgfplotsset{compat=1.18}
\usepackage[margin=1in]{geometry}
\usepackage{amsmath,amsfonts,amssymb,amsthm}
\usepackage[ruled,vlined,linesnumbered,commentsnumbered]{algorithm2e}
\SetKwInput{KwIn}{Input}
\SetKwInput{KwOut}{Output}
\SetKwInput{KwGlobal}{Global variable}
\SetKwInput{KwOracle}{Oracle}
\SetKwComment{Comment}{// }{}
\usepackage{graphicx}
\usepackage{enumerate}
\usepackage{bbm}
\usepackage{verbatim}
\usepackage{hyperref,color}
\usepackage[capitalize,nameinlink]{cleveref}
\usepackage[dvipsnames]{xcolor}
\hypersetup{
	colorlinks=true,
	pdfpagemode=UseNone,
	citecolor=OliveGreen,
	linkcolor=NavyBlue,
	urlcolor=Magenta,
	pdfstartview=FitW
}
\usepackage{appendix}
\usepackage{makecell}
\usepackage{tablefootnote}
\crefname{appsec}{Appendix}{Appendices}
\usepackage{tikz}
\usepackage{pgfplots}
\usepackage{xifthen}
\usepackage{subcaption}
\usepackage{hhline}
\usepackage{listings}

\theoremstyle{plain}
\newtheorem{theorem}{Theorem}[section]
\newtheorem{proposition}[theorem]{Proposition}
\newtheorem{lemma}[theorem]{Lemma}
\newtheorem{corollary}[theorem]{Corollary}
\newtheorem{conjecture}[theorem]{Conjecture}
\newtheorem{problem}[theorem]{Problem}
\newtheorem{claim}[theorem]{Claim}
\newtheorem{fact}[theorem]{Fact}

\theoremstyle{definition}
\newtheorem{definition}[theorem]{Definition}
\newtheorem{example}[theorem]{Example}
\newtheorem{observation}[theorem]{Observation}
\newtheorem{assumption}[theorem]{Assumption}
\newtheorem*{assumption*}{Assumption}
\newtheorem{condition}[theorem]{Condition}

\theoremstyle{remark}
\newtheorem{remark}[theorem]{Remark}

\crefname{lemma}{Lemma}{Lemmas}
\crefname{theorem}{Theorem}{Theorems}
\crefname{definition}{Definition}{Definitions}
\crefname{fact}{Fact}{Facts}
\crefname{claim}{Claim}{Claims}
\crefname{proposition}{Proposition}{Propositions}


\newcommand{\dif}{\,\mathrm{d}}
%\newcommand{\E}{\mathbb{E}}
%\newcommand{\Var}{\mathrm{Var}}
\newcommand{\Cov}{\mathrm{Cov}}
\newcommand{\MI}{\mathrm{I}}
%\newcommand{\Ent}{\mathrm{Ent}}
\newcommand{\ind}{\mathbbm{1}}
\newcommand{\allone}{\mathbf{1}}
\newcommand{\tr}{\mathrm{tr}}
\newcommand{\wt}{\mathrm{weight}}
\DeclareMathOperator*{\argmax}{arg\,max}
\newcommand{\norm}[1]{\left\lVert #1 \right\rVert}
\newcommand{\TV}[2]{\left\|#1 - #2\right\|_{\mathrm{TV}}}
\newcommand{\ceil}[1]{\left\lceil #1 \right\rceil}
\newcommand{\floor}[1]{\left\lfloor #1 \right\rfloor}
\newcommand{\ip}[2]{\left\langle #1 , #2 \right\rangle}
\newcommand{\grad}{\nabla}
\newcommand{\hessian}{\nabla^2}
\newcommand{\trans}{\intercal}
\newcommand{\diag}{\mathrm{diag}}
\newcommand{\poly}{\mathrm{poly}}
\newcommand{\dist}{\mathrm{dist}}
\newcommand{\sgn}{\mathrm{sgn}}
\newcommand{\fpras}{\mathsf{FPRAS}}
\newcommand{\fpaus}{\mathsf{FPAUS}}
\newcommand{\fptas}{\mathsf{FPTAS}}

\newcommand{\eps}{\varepsilon}

\newcommand{\N}{\mathbb{N}}
\newcommand{\Z}{\mathbb{Z}}
\newcommand{\Q}{\mathbb{Q}}
\newcommand{\R}{\mathbb{R}}
\newcommand{\C}{\mathbb{C}}

\renewcommand{\AA}{\mathcal{A}}
\newcommand{\BB}{\mathcal{B}}
\newcommand{\CC}{\mathcal{C}}
\newcommand{\DD}{\mathcal{D}}
\newcommand{\EE}{\mathcal{E}}
\newcommand{\FF}{\mathcal{F}}
\newcommand{\GG}{\mathcal{G}}
\newcommand{\HH}{\mathcal{H}}
\newcommand{\II}{\mathcal{I}}
\newcommand{\JJ}{\mathcal{J}}
\newcommand{\KK}{\mathcal{K}}
\newcommand{\LL}{\mathcal{L}}
\newcommand{\MM}{\mathcal{M}}
\newcommand{\NN}{\mathcal{N}}
\newcommand{\OO}{\mathcal{O}}
\newcommand{\PP}{\mathcal{P}}
\newcommand{\QQ}{\mathcal{Q}}
\newcommand{\RR}{\mathcal{R}}
\renewcommand{\SS}{\mathcal{S}}
\newcommand{\TT}{\mathcal{T}}
\newcommand{\UU}{\mathcal{U}}
\newcommand{\VV}{\mathcal{V}}
\newcommand{\WW}{\mathcal{W}}
\newcommand{\XX}{\mathcal{X}}
\newcommand{\YY}{\mathcal{Y}}
\newcommand{\ZZ}{\mathcal{Z}}

\newcommand{\KL}[2]{D_{\mathrm{KL}} \left( #1 \,\Vert\, #2 \right)}
\newcommand{\KLbig}[2]{D_{\mathrm{KL}} \left( #1 \,\big\Vert\, #2 \right)}
\newcommand{\KLBig}[2]{D_{\mathrm{KL}} \left( #1 \,\Big\Vert\, #2 \right)}

\newcommand{\SKL}[2]{D_{\mathrm{SKL}} \left( #1 \,\Vert\, #2 \right)}
\newcommand{\SKLbig}[2]{D_{\mathrm{SKL}} \left( #1 \,\big\Vert\, #2 \right)}
\newcommand{\SKLBig}[2]{D_{\mathrm{SKL}} \left( #1 \,\Big\Vert\, #2 \right)}

\newcommand{\qandq}{\quad\text{and}\quad}
\newcommand{\supp}{\mathrm{supp}}

\newcommand{\DKL}[2]{\-D_{\-{KL}}\left(#1 \parallel #2\right)}
\newcommand{\DTV}[2]{\-D_{\mathrm{TV}}\left({#1},{#2}\right)}
% \newcommand{\dist}{\mathrm{dist}}
\newcommand{\e}{\mathrm{e}}
\renewcommand{\epsilon}{\varepsilon}
\renewcommand{\emptyset}{\varnothing}
% \newcommand{\norm}[1]{\left\Vert#1\right\Vert}
\newcommand{\set}[1]{\left\{#1\right\}}
\newcommand{\tuple}[1]{\left(#1\right)} 
% \newcommand{\eps}{\varepsilon}
\newcommand{\inner}[2]{\left\langle #1,#2\right\rangle}
\newcommand{\tp}{\tuple}
\newcommand{\ol}{\overline}
\newcommand{\abs}[1]{\left\vert#1\right\vert}
\newcommand{\ctp}[1]{\left\lceil#1\right\rceil}
\newcommand{\ftp}[1]{\left\lfloor#1\right\rfloor}
\newcommand{\Ind}[1]{\-{Ind}_{#1}}
\newcommand{\todo}[1]{{\color{red} TODO: #1}}

\newcommand{\Cor}{\Psi^{\mathrm{sym}}}
\newcommand{\cor}{\Psi^{\-{cor}}}
\newcommand{\Inf}{\Psi}
% \newcommand{\diag}{\-{diag}}

\def\*#1{\boldsymbol{#1}} % Use \*A for \mathbf{A}
\def\+#1{\mathcal{#1}} % Use \+A for \mathcal{A}
\def\-#1{\mathrm{#1}} % Use \-A for \mathrm{A}
\def\=#1{\mathbb{#1}} % Use \^A for \mathbb{A}
\def\!#1{\mathfrak{#1}} % Use \!A for \mathfrak{A}

\newcommand*\diff{\mathop{}\!\mathrm{d}}

% probability
% \DeclareMathOperator*{\oPr}{\mathbf{Pr}}
\def\oPr{\mathop{\mathrm{Pr}}}
\renewcommand{\Pr}[2][]{ \ifthenelse{\isempty{#1}}
  {\oPr\left[#2\right]}
  {\oPr_{#1}\left[#2\right]} } % Use \Pr[a]{b} for \mathbf{Pr}_a[b], \Pr{b} for  \mathbf{Pr}[b]

% expectation: relies on the package "dsfont"
% \DeclareMathOperator*{\oE}{\mathbb{E}}
\def\oE{\mathop{\mathbb{E}}}
\newcommand{\E}[2][]{ \ifthenelse{\isempty{#1}}
  {\oE\left[#2\right]}
  {\oE_{#1}\left[#2\right]} }

% variance
%\DeclareMathOperator*{\oVar}{\mathbf{Var}}
\def\oVar{\mathrm{Var}}
\newcommand{\Var}[2][]{ \ifthenelse{\isempty{#1}}
  {\oVar\left[#2\right]}
  {\oVar_{#1}\left[#2\right]} }

% entropy
% \DeclareMathOperator*{\oEnt}{\mathbf{Ent}}
\def\oEnt{\mathrm{Ent}}
\newcommand{\Ent}[2][]{ \ifthenelse{\isempty{#1}}
  {\oEnt\left[#2\right]}
  {\oEnt_{#1}\left[#2\right]} }

\newcommand{\PhiEnt}[2][]{ \ifthenelse{\isempty{#1}}
  {\oEnt^\phi\left[#2\right]}
  {\oEnt^\phi_{#1}\left[#2\right]} }


\newenvironment{Exampleparagraph}[1][]
{
    \par 
    \noindent 
    \textbf{\color{blue} #1}
    \par 
    \vspace{0.5em} 
    \itshape 
}
{
    \par 
    \vspace{1em} 
}

%\input{local-macro}

\newcommand{\RED}[1]{\textcolor{red}{#1}}

\newcommand{\bern}[1]{\mathrm{Bern}\left(#1\right)}
\newcommand{\pll}{\parallel}



\title{CS 7545 F26, Homework 2}
\date{Due on October 1, 2026}
\author{{\textbf{Zejia Chen}} \\ \textbf{GTID: 904178159}}

\pgfplotsset{compat=1.18} 
\begin{document}

\maketitle

\section*{Exercise 1}
\textbf{(1)} 
Suppose $h^\star = (b_1^\star,\cdots,b_d^\star)$ and let $j_i^\star \in \arg\min_{j\in[R]} \abs{\frac{2j-1}{2R} - b_i^\star}$.
Let $h' = \tp{\frac{2j_1^\star - 1}{2R},\cdots,\frac{2j_d^\star -1 }{2R}}$.
Define 
$$
  \ell_i = \min\tp{b_i^\star,\frac{2j_i^\star-1}{2R}}, r_i = \max\tp{b_i^\star,\frac{2j_i^\star-1}{2R}}.
$$
Then we have that
\begin{align*}
  L(h_R^\star) \leq L(h') &= \Pr[(x,y)]{\ind\left[\forall i,x_i\leq \frac{2j_i^\star - 1}{2R}\right] \neq \frac{y+1}{2}} 
  \\&\leq \Pr[(x,y)]{\tp{\ind[\forall i, x_i\leq b_i^\star] \neq \frac{y+1}2 }\lor \tp{\exists i, \ell_i<x_i\leq r_i}}
  \\&\leq L(h^\star) + \sum_{i=1}^d \Pr[x]{\ell_i<x_i\leq r_i}
  \\&= L(h^\star) + \sum_{i=1}^d r_i-\ell_i
  \leq L(h^\star) + \frac{d}{2R},
\end{align*}
where the second inequality follows from the union bound and the last inequality follows from $\min_{j\in[R]} \abs{\frac{2j-1}{2R} - b_i^\star} \leq \frac 1{2R}$ for $i\in [d]$.
\\\textbf{(2)}
By Hoeffding's bound, for any fixed $h\in H_R$, we have that 
$$
  \Pr{\hat L_S(h) - L(h) \geq \eps } \leq \exp(-2n\eps^2), \quad
  \Pr{L(h) - \hat L_S(h) \geq \eps } \leq \exp(-2n\eps^2).
$$
Therefore, by union bound, with probability at least $1 - 2R^{d}\exp(-2n\eps^2)$, we have that
$$
  \forall h\in H_R, L(h) - \hat L_S(h) + \hat L_S(h_R^\star) - L(h_R^\star)\leq 2\eps
  % \implies \forall h\in H, L(h) - L(h_R^\star) \leq \eps, 
$$
Since $\hat L_S(h_R^\star) \geq \hat L_S(\hat h_{S,R})$, we have that 
\begin{align}
  \Pr{L(\hat h_{S,R}) - L(h_{R}^\star)\leq 2\eps}
  &\geq \Pr{L(\hat h_{S,R}) - \hat L_S(\hat h_{S,R}) + \hat L_S(h_R^\star) - L(h_{R}^\star)\leq 2\eps} \nonumber
  \\&\geq \Pr{\forall h\in H_R, L(h) - \hat L_S(h) + \hat L_S(h_R^\star) - L(h_{R}^\star)\leq 2\eps} \nonumber
  \\&\geq 1 - 2R^{d}\exp(-2n\eps^2). \label{eq:E1_P2}
\end{align}
Plugging $\eps = \sqrt{\frac{d\ln(R)+\ln(2/\delta)}{2n}}$ into \Cref{eq:E1_P2}, we complete the proof.
\\\textbf{(3)}
By Part \textbf{(1)} and Part \textbf{(2)}, take $R= dn $, with probability at least $1 - \delta$, we have that 
$$
  L(\hat h_{S,R}) - L(h^\star) \leq \frac{d}{2R} + 2\sqrt{\frac{d\ln(R) + \ln(2/\delta)}{2n}} \lesssim \frac 1n + \sqrt{\frac{d\ln(dn)}{n}} = O\tp{\sqrt{\frac{d\ln (dn)}{n}}}.
$$

\section*{Exercise 2}
\textbf{(1)} The VC dimension is $2$.
For $S\subseteq \+X$ with $|S| = 3$, we have that $(1,1,1)\notin H\mid_S$ since the length of the interval in $H$ is only $1.5$ and $\sup S - \inf S \geq 2$.
For $S = \{1,2\}$, it is clear that $H\mid_S = \{(-1,-1),(-1,1),(1,-1),(1,1)\}$.
\\\textbf{(2)}$\Pi_H(n) = 2n$.
Since the VC dimension is $2$, $\Pi_H(n) = 2^n = 2n$ for $n\leq 2$.
For $n\geq 3$, since the length of the interval in $H$ is only $1.5$, it can only cover at most $2$ points and these two points must be consecutive.
Therefore, for any $S \subset \Z$ with $|S| = n$, we have that  
$$
  \abs{H|_S} = 1 + n + \sum_{x\in S}\ind [x+1\in S] \leq 1 +n + n-1 = 2n,
$$
with equality achieved by $S = [n]$.
Therefore, $\Pi_H(n) = \max_{|S|=n} \abs{H|_S} = 2n$.
\\\textbf{(3)}
The upper bound of the Sauer-Shelah Lemma is $\frac{n^2+n+2}{2}$, which is quadratic in $n$, whereas $\Pi_H(n)$ is linear in $n$.
The underlying reason is that the general upper bound misses the combinatorial structure of $H$.
Here if two points in $S$ are labeled $+1$, these two points must be consecutive, which is not generally true for hypothesis class with VC dimension $2$.
This property restricts $\Pi_H(n)$ to $2n$.
\section*{Exercise 3}
\textbf{(1)}
We have that
\begin{align*}
  \Pi_{H_1\cup H_2}(n) &= \sup_{|S|=n}\abs{(H_1\cup H_2)|_S} \leq \sup_{|S|=n}\tp{\abs{H_1|_S} + \abs{H_2|_S}}
  \\&\leq \sup_{|S_1|=n}\abs{H_1|_{S_1}} + \sup_{|S_2|=n}\abs{H_2|_{S_2}} = \Pi_{H_1}(n) + \Pi_{H_2}(n)
\end{align*}
\textbf{(2)}
By Part \textbf{(1)} and the Sauer-Shelah Lemma, we have that for any set $S$ with $|S| = d_1+d_2+2$,
\begin{align*}
  \Pi_{H_1\cup H_2}(d_1+d_2+2) &\leq \Pi_{H_1}(d_1+d_2+2) + \Pi_{H_2}(d_1+d_2+2)
  \\&\leq \sum_{i=0}^{d_1} \binom{d_1+d_2+2}{i} + \sum_{i=0}^{d_2} \binom{d_1+d_2+2}{(d_1+d_2+2)-i}
  \\&= 2^{d_1+d_2+2}-\binom{d_1+d_2+2}{d_1+1}
  < 2^{d_1+d_2+2}.
\end{align*}
Therefore, $d(H_1\cup H_2)\leq d_1+d_2+1$.
\\\textbf{(3)}
For any set $S\subset \Z$ with $|S| = d_1$, we have that $\abs{H_1|_S} = |\{T\subseteq S\mid |T|\leq d_1\}| = 2^{d_1}$.
However, for any $S\subset \Z$ with $|S| = d_1 + 1$, $(1,1,\cdots,1)\notin H_1|_S$.
Therefore, $d(H_1) = d_1$.
In the same way, $d(H_2) = d_2$.

Fix $S\subset Z$ with $|S| = d_1 + d_2 + 1$.
For any vector $v$ in $\{\pm 1\}^{|S|}$, either $v$ contains at most $d_1$ entries equal to $+1$ or at most $d_2$ entries equal to $-1$.
In the former case, there exists some $h\in H_1$ such that $(h(x_1),\cdots,h(x_{|S|})) = v$.
In the other case, $v$ can be achieved by some $h\in H_2$.
Therefore, $\abs{(H_1\cup H_2)|_S} = 2^{d_1+d_2+1}$.
By Part \textbf{(2)}, we conclude that $d(H_1\cup H_2)  = d_1+d_2+1$.
\section*{Exercise 4}
\textbf{(1)}
By the union bound and Hoeffding's bound, we have that
\begin{equation}\label{eq:E4_P1}
  \Pr{\exists h\in H_k, L(h) - \hat L_S(h) \geq \eps} \leq \sum_{h\in H_k} \cdot \Pr{L(h) - \hat L_S(h) \geq \eps} \leq \exp(k-2n\eps^2).
\end{equation}
Plugging $\eps = \sqrt{\frac{k+\ln(1/\delta)}{2n}}$ into \Cref{eq:E4_P1}, we have that
\begin{align*}
  \Pr{L(\hat h_{S,k}) \leq \hat L_S(\hat h_{S,k}) + \sqrt{\frac {k+\ln (1/\delta)}{2n}}} \geq \Pr{\forall h\in H_k, L(h) - \hat L_S(h) \leq \sqrt{\frac {k+\ln (1/\delta)}{2n}}} \geq 1 - \delta.
\end{align*}
\textbf{(2)}
Since for any hypothesis $h$, $n\hat L_S(h)\in \Z_{\geq 0}$, we have that for $k < k_S$, $\min_{h\in H_{k+1}} \hat L_S(h) \leq \min_{h\in H_{k}} \hat L_S(h) - \frac 1n$.
The loss function satisfies $\hat L_S(h)\in [0,1]$ for hypothesis $h$.
Therefore, we conclude that
$$
  0\leq \min_{h\in H_{k_S}} \hat L_S(h) \leq \min_{h\in H_{k_S-1}} \hat L_S(h)-\frac 1n \leq \cdots \leq \min_{h\in H_{0}} \hat L_S(h)-\frac {k_S}n \leq 1-\frac {k_S}n ,
$$
which shows $k_S\leq n$.
\\\textbf{(3)}
By union bound and Part \textbf{(1)}, for $\delta' > 0$, with probability at least $1 - (n+1)\delta'$, 
$$
  \forall k\in \{0\}\cup [n], \quad L(\hat h_{S,k}) \leq \hat L_S(\hat h_{S,k}) + \sqrt{\frac {k+\ln (1/\delta')}{2n}}.
$$
By Part \textbf{(2)} and $\hat h_{S,\-{iter}} = \hat h_{S,k_S}$, and take $\delta = \delta'(n+1)$, we conclude that
\begin{align*}
  &\Pr{L(\hat h_{S,k_S}) \leq \hat L_S(\hat h_{S,k_S}) + \sqrt{\frac {k_S+\ln(n+1)+\ln (1/\delta)}{2n}}} 
  \\\geq &\Pr{\forall k\in \{0\}\cup [n], L(\hat h_{S,k}) \leq \hat L_S(\hat h_{S,k}) + \sqrt{\frac {k+\ln(n+1)+\ln (1/\delta)}{2n}}} \geq 1 - \delta.
\end{align*}

\section*{Exercise 5}
\textbf{(1)}
For each point $x$, sample its label from $\{\pm 1\}$ uniformly at random and the sampling is independent of the algorithm.
Let $\+X = \{x_1,\cdots,x_d\}$ be the data points and $\+Y$ be the corresponding sampled labels.
Fix the sample set $S$, for any output $h$ of $\+A(S)$, since $\+X$ is shattered by $H$, we have that for any $x_i\in \+X$,
$$
  x_i\notin S\implies y_i = h(x_i)\ \text{w.p.}\ \frac 12 .
$$ 
Therefore, any excluded point will contribute error $\frac 12$.
From this, we conclude the expectation of the loss function of non-single points:
\begin{align*}
  \E[\+Y]{\E[S,\+A]{L_{\text{non-single}}(\+A(S))}} &\geq \frac 12(d-1)\cdot \frac{b}{d-1}\cdot \tp{1 - \frac{b}{d-1}}^n
  \\&\geq  \frac 12 b\cdot \tp{1 - \frac{nb}{d-1}} ,
\end{align*}
where the last inequality follows from the Bernoulli's inequality.
Choosing $b = \frac{d-1}{2n}$, we get $\E[\+Y]{\E[S,\+A]{L_{\text{non-single}}(\+A(S))}} \geq \frac{d-1}{8n}$.
\\\textbf{(2)}
By Part \textbf{(1)}, we have that
$$
  \max_{\+Y\in \{\pm 1\}^{\+X}} \E[S,\+A]{L_{\text{non-single}}(\+A(S))}
  \geq \E[\+Y]{\E[S,\+A]{L_{\text{non-single}}(\+A(S))}} = \frac b2\tp{1 - \frac{nb}{d-1}} .
$$
Fix $\+Y\in \arg \max_{\+Y\in \{\pm 1\}^{\+X}} \E[S,\+A]{L_{\text{non-single}}(\+A(S))}$.
We get the distribution $\+D$ from Part \textbf{(1)}, which is realizable by $H$.
Let $Z = L_{\text{non-single}}(\+A(S))$.
Since $\+D(\text{non-single point}) = b$, we have that $Z \leq b$ and $b-Z$ is nonnegative.
Take $b = \frac{d-1}{2n}$, then by Part \textbf{(1)} and the choice of $\+Y$, $\E{Z} \geq b/4$.
By Markov's inequality, we have that
$$
  \Pr{Z\leq \frac{b}{8}} = \Pr{b-Z\geq \frac{7b}{8}} \leq \frac{b-\E{Z}}{7b/8} \leq \frac{6}{7}.
$$
Therefore, since $d\geq 2$,
$$  
  \Pr{L(\+A(S))< \frac 1{32}\frac dn} \leq \Pr{Z< \frac 1{32}\frac dn}
  \leq \Pr{Z \leq \frac 1{16}\frac {d-1}n} \leq \frac 67.
$$
\textbf{(3)}
Let these two points be $x_1,x_2$ and $\+D(x_1) = 1-b$.
Let $S = \{s_1,\cdots,s_n\}$.
Since $s_1,\cdots,s_n$ are sampled i.i.d., we have that 
$$
  \Pr{\forall i\in[n], s_i = x_1} = \prod_{i=1}^n \Pr{s_i=x_1} = \+D(x_1)^n = (1-b)^n .
$$
As in Part \textbf{(1)}, we sample the labels $\+Y$ of data points from $\{\pm 1\}$ uniformly at random.
If $x_2\notin S$, then $\E{\ind[\+A(S)(x_2)\neq y_2]} = \frac 12$.
Therefore, there exists $y_2' \in \{\pm 1\}$ such that $\Pr[\+A]{\+A(S)(x_2)\neq y_2'}\geq \frac 12$ (note that here the randomness comes from the algorithm).
Hence, we have that
$$
  \Pr[S,\+A]{\+A(S)(x_2)\neq y_2'} \geq \Pr[\+A]{\+A(\{x_1,x_1,\cdots,x_1\})(x_2)\neq y_2'}\Pr{x_2\notin S} \geq (1-b)^n/2 .
$$
\textbf{(4)}
When the other point is incorrectly classified, the loss is at least $b$ ($\+D(x_2) = b$).
Hence, with probability at least $(1-b)^n/2$, $L(\+A(S))\geq b$.

Take $b = \frac{\ln(1/(2\delta))}{2n}$.
Since $\delta\in [e^{-n},1/2]$, $b\in [0,\frac 12]$.
Then $(1-b)^n/2 \geq \exp(-2bn)/2$.
Therefore,
$$
  \Pr{L(\+A(S)) \geq \frac{\ln(1/(2\delta))}{2n}} \geq \frac{(1-b)^n}2\geq \frac{e^{-2bn}}2 = \delta.
$$
\end{document}